\uuid{fhUI}
\titre{Équation droite/plan}
\auteur{bodin ; hanani ; mzari}
\organisation{exo7}
\licence{CC BY-NC-SA 4.0}
\datecreate{2020-01-20}
\niveau{L1}
\module{Géométrie}
\chapitre{Géométrie de l'espace}
\sousChapitre{Plans}
\difficulte{3}
\theme{Plans -- Droites, équation droite, plan}
\exo7qcmid{lille-1.geometrie-espace.016}

\contenu{
\texte{
Pour ces questions, l'espace est muni d'un repère orthonormé direct  $(O, \vec{i}, \vec{j}, \vec{k})$.
}

\begin{question}
\qclassification{141.03,141.04}
\qtags{motcle=équation droite/plan}
Soit $D$ la droite d'équations :  $ \left\{\begin{array}{ccl}x+z&=&1\\x-y&=&-1 \end{array}\right.$ et $P$ le plan contenant $D$ et perpendiculaire au plan $Q$ d'équation : $x-z+3=0$.  Une équation cartésienne de $P$ est :
\begin{answers}

    \good{$x+z=1$}

      \bad{$x+y=0$}

     \bad{$y+z=1$}

     \bad{$x-y=-1$}

\end{answers}
\begin{explanations}
$\vec{u}(1,0,-1)$ est un vecteur normal à $ Q$ qui n'appartient pas au  plan vectoriel $x-y=0$. Donc $P$ est différent du plan d'équation : $x-y=-1$ et donc  une équation cartésienne de $P$ est  de la forme : $(x+z-1) + \alpha(x-y+1)=0, \, \alpha \in \Rr$. On calcule $\alpha$ de sorte que $\vec{u}(1,0,-1)$ soit un vecteur de  $P$.
\end{explanations}
\end{question}
}
